%% kapa: exam=1 round=24 year=2013 period=1 form=A subject=economics subject_ko=경제원론 seq=8 number=48
%% source: Q-Net 감정평가사 기출문제 (cst00309) artlSeq=5073802 file=extracted/1008/5073802/2013년 제24회 감정평가사 1차시험 기출문제/1교시(A형).hwp
%% conversion: kordoc 4.24.51 -> tools/gen1.py (auto); review=proofread-sonnet; flags=none
\begin{question}{48}
현재 소비자 甲은 주어진 소득을 모두 사용하여 가격이 1,000원인 A재 10단위와 500원인 B재 15단위의 조합을 소비하려고 한다. 이때의 한계대체율$(MRS_{ AB } = \frac{ M U_{ A } } { M U_{ B } } )$이 1.5라면 효용극대화를 위한 甲의 선택으로 옳은 것은? (단, 소비자 甲의 무차별곡선은 우하향하고 원점에 대하여 볼록하며, $M U_{ i }$는 $i$재의 한계효용이다.)
\begin{choices}
  \item A재 10단위와 B재 15단위를 소비한다.
  \item A재와 B재 소비를 모두 증가시켜야 한다.
  \item A재와 B재 소비를 모두 감소시켜야 한다.
  \item A재 소비를 증가시키고, B재 소비를 감소시켜야 한다.
  \item A재 소비를 감소시키고, B재 소비를 증가시켜야 한다.
\end{choices}
\end{question}
\answer{5}
